\subsubsection{A residue-class multiplicity bound for the Fitting error}
\label{sec:adriazola-fitting-multiplicity}
The factorial correction in the preceding generalization of
Adriazola's Fitting argument \cite{adriazola-reconstruction}
can be localized to repeated residual simple representations.
Keep its determinant-tuple and off-target-kernel assumptions,
and assume individual residual surjectivity as there.
Write \(v(\pi)=1\), and first let the residue field have finite
cardinality \(q\). For \(t\ge0\) integral put
\[
 w_q(t)=\sum_{h\ge1}\left\lfloor t/q^h\right\rfloor .
\]
Partition the diagonal representations \(j\) by the kernels of
\(R\to M_{n_j}(k)\) modulo \(\pi\). For a class \(c\) write \(m_c\)
for its multiplicity and \(n_c\) for its common matrix size.

\begin{theorem}[Multiplicity-sensitive Fitting bound]
\label{thm:fitting-multiplicity}
With the notation of the preceding Fitting estimate, one has
\[
 \delta_{x_0}\le\sum_j L_j+
                   b\sum_c n_c w_q(m_c).                 \tag{MF1}
\]
In particular the exact bound holds if every \(m_c<q\), irrespective
of the total number of distinct residual blocks or their dimensions.
No determinant law after scalar extension is needed. The statement
holds also for an equal-characteristic DVR. If the residue field
is infinite, the error is zero.
\end{theorem}
\begin{proof}
We first prove the interpolation estimate used below. Choose representatives
\(u_0,\ldots,u_{q-1}\) of the residue field, with \(u_0=0\).
For each nonnegative integer \(i\) write its base-\(q\) expansion
\(i=\sum_h i_hq^h\) and set \(a_i=\sum_h u_{i_h}\pi^h\).
For distinct \(i,j\), their first unequal digit gives
\[
 v(a_i-a_j)=\max\{h:i\equiv j\bmod q^h\}.
\]
In particular the product \(d_t=\prod_{0\le i<t}(a_t-a_i)\)
has valuation \(w_q(t)\).
For every \(x\in O\), at least \(\lfloor t/q^h\rfloor\) of the first
\(t\) points coincide with \(x\) modulo \(\pi^h\), by the
distribution of the first \(h\) digits. Hence
\[
 N_t(X)=d_t^{-1}\prod_{0\le i<t}(X-a_i)
\]
takes integral values on \(O\), is zero at \(a_i\) for \(i<t\),
and equals one at \(a_t\). Include \(N_0=1\).
For a polynomial \(F\) in finitely many variables of total degree
at most \(D\), expand \(F=\sum_{|\alpha|\le D}
c_\alpha\prod_i N_{\alpha_i}(X_i)\). At the points
\((a_{\alpha_1},\ldots,a_{\alpha_t})\), ordered by total degree,
the evaluation matrix is unit lower triangular with integral
entries: a nonzero entry requires coordinatewise inequality.
Thus if all these values have valuation at least \(B\), all
\(c_\alpha\) have valuation at least \(B\).
If \(F\) is homogeneous of degree \(D\) and the coefficient of
\(X^\alpha\), \(|\alpha|=D\), is a unit, only this Newton basis
element contributes to that ordinary coefficient. Consequently
\(v(c_\alpha)=\sum_i w_q(\alpha_i)\), and some evaluation of \(F\)
has valuation at most that number. This proof uses only finite
polynomial identities, not completeness.

Return to the block data. We may assume every \(L_j\) is finite.
As in the preceding Fitting proof, choose bases of the saturated
images \(W_j\), and let
\[
 \Phi(C)=\prod_j\det(\rho_j(C)|_{W_j})
\]
for \(C\in M_{b\times a}(R)\). We now specify its coefficient
variables. The kernels of the residual surjections are maximal
two-sided ideals in \(R\), because each quotient is a simple
matrix algebra. Distinct such kernels are comaximal. The
Chinese remainder map to the product of the distinct quotients
is surjective: for a fixed quotient choose, in each other
kernel, an element congruent to one in the fixed quotient,
and multiply them. The resulting element vanishes in all
other quotients and equals one in the fixed quotient. Summing
lifts multiplied by these selectors proves surjectivity.
Within one class the identifications with full matrix algebras
differ by a \(k\)-algebra automorphism, hence by conjugation:
this follows, for example, by transporting the matrix units
and choosing a basis of their one-dimensional image summands.
In particular their sizes agree.

Choose lifts of the matrix-unit basis for the product of the
distinct quotients. Parameterize each entry of \(C\) by
\(O\)-linear combinations of those lifts, with independent
variables. Then \(\Phi\) is homogeneous of degree
\(D_0=b\sum_j n_j\). Modulo \(\pi\), make the just-described
basis changes in the source and target of each block.
Up to nonzero determinant factors, the residual polynomial for
class \(c\) is a product of \(m_c\) determinants
\(\det(X_c B_j)\); here \(X_c\) is the full
\(bn_c\)-by-\(an_c\) variable matrix for that quotient and
\(B_j\) is the full-column-rank matrix of the transformed
saturated subspace. Each determinant is nonzero as a polynomial:
\(X_c\) can be chosen a left inverse of \(B_j\).
The product over all classes is nonzero in the polynomial ring.
It is separately homogeneous of degree \(m_c\) in the variables
of each of the \(bn_c\) output rows of \(X_c\).
Choose any nonzero monomial in this residual polynomial. Its
multi-index \(\alpha\) therefore satisfies
\[
 \sum_i w_q(\alpha_i)
 \le\sum_c bn_c w_q(m_c)
\]
by the row sums and superadditivity of the floor functions.
The corresponding ordinary coefficient of the original
homogeneous \(\Phi\) is a unit. Apply the interpolation lemma
to obtain \(C\) with finite \(v(\Phi(C))\) at most this bound.
As in the preceding argument, the determinant tuple for \(CP\)
lies in the original algebra and vanishes off \(x_0\); its
valuation at \(x_0\) is exactly \(\sum_jL_j+v(\Phi(C))\).
This proves (MF1). For an infinite residue field, a nonzero
polynomial has a nonzero residue value, which gives error zero.
\end{proof}

For mixed characteristic \(\ell\), normalize (MF1) by \(e=v(\ell)\).
The correction is \(b\sum_c n_c w_q(m_c)/e\), at most
\(b\sum_c n_c v_\ell(m_c!)\), since \(q\) is a power of \(\ell\)
and \(e\ge1\). This is no larger than the previous
\(v_\ell(D_0!)\), by superadditivity. If the multiplicities
are bounded by \(m\), the error vanishes for \(\ell>m\).
For example a long collection of pairwise inequivalent residual
blocks incurs no error even over \(\mathbf F_2\); a collection of
repeated scalar blocks with multiplicity three costs at most
\(b\,w_q(3)\) per class, zero over \(\mathbf F_4\).
For an arithmetic application the classes and their multiplicities
must be those of the actual residual representations, and the
original determinant and off-target kernel conditions still apply.
No global weighted Selmer estimate follows from the algebra alone.
